\chapter{La fonction de gestion des risques}\label{ch:fm:the-risk-management-function}

L'examinateur nommé dans la faillite de Lehman Brothers a indiqué en mars 2010 qu'entre décembre 2006 et décembre 2007 la société a relevé
sa limite d'\omterm{def:fm:the-head-of-desks-job:appetite}{appétit pour le risque} à l'échelle de la société à trois reprises, de 2.3~milliards de dollars à 4.0~milliards de dollars. La dernière hausse a été approuvée en janvier 2008 et datée rétroactivement au
décembre 2007, ce qui a fait disparaître du registre les dépassements des semaines intermédiaires~; les conseillers de l'examinateur ont calculé que la méthode utilisée
l'année précédente aurait donné une limite d'environ 2.5~milliards de dollars. Les tests de résistance de la société ont laissé de côté ses investissements en immobilier commercial et en capital-investissement
ainsi que, pendant un temps, ses prêts à effet de levier. Les chiffres étaient là. Ce qui a fait défaut, c'était une fonction incapable de leur donner du poids.

\section{Indépendance~: à qui rend compte la fonction de gestion des risques}

La fonction de gestion des risques d'une société de trading mesure les risques que prennent ses activités, fixe et contrôle les limites du chapitre~7, et formule une objection lorsque les deux divergent.
C'est la deuxième des trois lignes de défense du livre~6~: les activités sont responsables de leurs risques, la fonction de gestion des risques les conteste, et l'audit interne vérifie
les deux. Sa valeur tient entièrement à l'objection, et une objection ne vaut que ce que vaut la ligne hiérarchique qui la soutient.

\begin{definition}[Directeur des risques, comité des risques]\label{def:fm:the-risk-management-function:cro}
Le \emph{directeur des risques}\index{directeur des risques} (chief risk officer) est le dirigeant responsable de la fonction de gestion des risques~: de la fixation des limites de risque pour l'ensemble de la société
et du suivi de leur respect, des politiques et des contrôles qui mettent en œuvre l'\omterm{def:fm:the-head-of-desks-job:appetite}{appétit pour le risque}, et de la remontée des dépassements, des défaillances
et des risques émergents. Le \emph{comité des risques}\index{comité des risques} (risk committee) est l'instance, du conseil d'administration ou de la direction générale, qui approuve les politiques de risque
et les limites, reçoit les rapports du directeur des risques et statue sur les dérogations et sur les hausses que les politiques lui réservent.
\end{definition}

Trois éléments rendent le \omterm{def:fm:the-risk-management-function:cro}{directeur des risques} indépendant dans les faits~: une ligne hiérarchique qui ne passe pas par les activités dont le risque est mesuré,
une rémunération qui ne dépend pas de leurs revenus, et un lien direct avec le \omterm{def:fm:the-risk-management-function:cro}{comité des risques} que nul intermédiaire ne peut filtrer. La première est visible sur un
organigramme (la \cref{fig:fm:the-risk-management-function:org})~; les autres ne le sont pas, et comptent tout autant. Le rapport de l'examinateur cite des articles de presse
publiés avant la faillite, selon lesquels la société avait écarté son \omterm{def:fm:the-risk-management-function:cro}{directeur des risques} en 2007 en raison de l'opposition à son
accumulation croissante de placements illiquides~; le \omterm{def:fm:the-risk-management-function:cro}{directeur des risques} avait aussi plaidé, fin 2006, pour une hausse de limite plus faible que celle qui a été adoptée.

\begin{omfigure}
\begin{tikzpicture}[font=\footnotesize,>=stealth,box/.style={draw,rounded corners,align=center,fill=omDef!10,minimum height=0.75cm,text width=2.7cm}]
\node[box,fill=gray!15] (board) at (0,3.3) {conseil};
\node[box,fill=omProp!15] (rc) at (-3.4,2.1) {comité des risques};
\node[box,fill=gray!15] (ceo) at (3.4,2.1) {directeur général};
\node[box,fill=omProp!15] (cro) at (-3.4,0.6) {directeur des risques};
\node[box] (bus) at (3.4,0.6) {responsable des marchés};
\node[box] (mr) at (-5.0,-0.9) {risque de marché et de crédit};
\node[box] (mv) at (-1.9,-0.9) {validation des modèles};
\node[box] (d1) at (1.9,-0.9) {desks};
\node[box] (d2) at (5.0,-0.9) {support au trading};
\draw[->] (board) -- (rc); \draw[->] (board) -- (ceo);
\draw[->,very thick,omProp] (rc) -- (cro); \draw[->,very thick,omProp] (ceo) -- (cro);
\draw[->] (ceo) -- (bus); \draw[->] (cro) -- (mr); \draw[->] (cro) -- (mv); \draw[->] (bus) -- (d1); \draw[->] (bus) -- (d2);
\draw[<->,dashed,gray] (cro) -- node[above,font=\scriptsize,text=black] {remise en cause} (bus);
\end{tikzpicture}

\omcaption{Une fonction de gestion des risques indépendante sur un organigramme~: le \omterm{def:fm:the-risk-management-function:cro}{directeur des risques} rend compte à la fois au \omterm{def:fm:the-risk-management-function:cro}{comité des risques} et au directeur général
(traits épais), et non au responsable des activités dont le risque est mesuré, et les conteste latéralement. Schéma simplifié.}\label{fig:fm:the-risk-management-function:org}
\end{omfigure}

\begin{dated}{2026-09}{Le directeur des risques dans les règles bancaires américaines}\label{dat:fm:the-risk-management-function:rules}
En vertu de la Regulation YY de la Federal Reserve (12 CFR 252.33), une société holding bancaire américaine dont les actifs consolidés totaux moyens atteignent 100~milliards de dollars ou plus
doit maintenir un \omterm{def:fm:the-risk-management-function:cro}{comité des risques} au sein de son conseil d'administration~: un comité indépendant doté d'une charte écrite, présidé par un administrateur indépendant, se réunissant au moins
chaque trimestre et consignant ses décisions. Elle doit nommer un \omterm{def:fm:the-risk-management-function:cro}{directeur des risques} expérimenté dans la gestion des risques des grands établissements financiers complexes, responsable
de fixer les limites de risque à l'échelle du groupe et d'en surveiller le respect~; sa rémunération doit être cohérente avec une évaluation objective des
risques pris, et il doit rendre compte directement au \omterm{def:fm:the-risk-management-function:cro}{comité des risques} et au directeur général. D'autres juridictions et des établissements plus petits fixent leurs propres
règles~; une société de trading pour compte propre peut n'être soumise à aucune exigence réglementaire.
\end{dated}

Les superviseurs qui ont comparé les grandes firmes après les premiers mois de la crise de 2007 ont constaté que la différence ne tenait pas à la seule indépendance. Le Senior
Supervisors Group a indiqué en mars 2008 que, dans les firmes qui ont évité des pertes importantes, la fonction de gestion des risques disposait d'indépendance et d'autorité, mais aussi
une interaction directe considérable avec les dirigeants des activités~; que certaines firmes ont fait remonter leurs inquiétudes sur les risques émergents dès l'été
2006, gagnant jusqu'à un an~; et que, dans d'autres, les structures hiérarchiques filtraient ce qui parvenait à la direction générale. Une fonction de gestion des risques indépendante
et distante voit les chiffres tardivement et est peu écoutée.

\section{Comités et outils du directeur des risques}

Les outils de la fonction de gestion des risques sont peu nombreux et grossiers. Elle propose l'\omterm{def:fm:the-head-of-desks-job:appetite}{appétit pour le risque} (chapitre~6) et le \omterm{def:fm:limits-and-capital-allocation:framework}{dispositif de limites} (chapitre~7)~; elle mesure la valeur en
risque et conduit des tests de résistance (livre~6)~; son équipe de validation des modèles approuve les modèles qui produisent ces chiffres~; elle valide les nouveaux produits et les nouvelles
stratégies~; elle décide, ou recommande, ce qui se passe lorsqu'une limite est dépassée. Chaque outil ne fonctionne que si une transaction ne peut pas avoir lieu sans lui.

\begin{definition}[Principe des quatre yeux]\label{def:fm:the-risk-management-function:foureyes}
Le \emph{principe des quatre yeux}\index{principe des quatre yeux} exige qu'une décision ou une action d'un type défini (une hausse de limite, une modification de modèle, un
nouveau produit, un paiement) soit approuvée par une seconde personne, indépendante de celle qui la propose ou l'exécute, avant de prendre effet.
\end{definition}

Ce principe est la forme opérationnelle de la séparation des tâches (livre~6)~: le trader qui veut une limite plus grande ne l'accorde pas, et le quant qui
modifie un modèle de risque ne valide pas lui-même la modification. Il échoue lorsque la seconde paire d'yeux appartient à quelqu'un dont les intérêts sont ceux de la première personne~:
un responsable de desk qui approuve les hausses pour son propre desk applique quatre yeux sur le papier et deux en réalité.

\begin{method}[Faire vivre un comité des risques]\label{met:fm:the-risk-management-function:committee}
\begin{enumerate}
\item Rédiger sa charte~: membres, quorum, ce qu'il est seul à pouvoir approuver (hausses au-delà d'un seuil, nouveaux produits, modifications de modèles qui font varier une mesure de
limite), et la fréquence de ses réunions.
\item Lui fournir des rapports réguliers émanant de la fonction de gestion des risques, et non des activités~: taux d'utilisation des limites, dépassements et leur résolution, résultats des tests de
résistance, dérogations en cours.
\item Consigner chaque décision avec sa motivation et son échéance~; réexaminer à chaque réunion les dérogations arrivées à échéance.
\item Accorder au \omterm{def:fm:the-risk-management-function:cro}{directeur des risques} le droit de saisir directement le comité du conseil d'administration, et s'en servir.
\end{enumerate}
\end{method}

\section{Escalade~: dépassements, excès et dérogations}

Le chapitre~7 a défini la notion de \omterm{def:fm:limits-and-capital-allocation:breach}{dépassement de limite} et la différence entre \omterm{def:fm:limits-and-capital-allocation:breach}{limites fermes} et \omterm{def:fm:limits-and-capital-allocation:breach}{limites indicatives}. Ce qui suit est la \omterm{def:fm:the-risk-management-function:escalation}{procédure d'escalade}.

\begin{definition}[Procédure d'escalade, dérogation de risque]\label{def:fm:the-risk-management-function:escalation}
Une \emph{procédure d'escalade}\index{procédure d'escalade} précise, pour chaque gravité de \omterm{def:fm:limits-and-capital-allocation:breach}{dépassement de limite} ou de défaillance de contrôle, qui doit être informé, dans
quel délai, qui peut approuver une augmentation temporaire ou une autre dérogation, et ce qu'il faut faire si aucune approbation n'est donnée. Une
\emph{dérogation de risque}\index{dérogation de risque} est un écart documenté, approuvé et limité dans le temps par rapport à une limite ou à une politique, avec un responsable et une date d'échéance.
\end{definition}

Un dépassement est classé selon sa taille et sa durée. Les règles du chapitre (le \cref{lst:fm:the-risk-management-function:rules}) classent un dépassement de gravité~1 si
il dépasse la \omterm{def:fm:limits-and-capital-allocation:breach}{limite ferme} de 5~\% au maximum pendant deux jours au plus, de gravité~3 s'il la dépasse de plus de 20~\% ou dure plus de cinq jours, et de
gravité~2 entre les deux. Le dépassement de gravité~1 est signalé au responsable de desk le jour même, celui de gravité~2 au \omterm{def:fm:the-risk-management-function:cro}{directeur des risques} dans la journée, et celui de gravité~3 au
\omterm{def:fm:the-risk-management-function:cro}{comité des risques} dans les deux jours. Il existe trois voies honnêtes pour sortir d'un dépassement~: réduire la position, faire approuver une hausse temporaire par
une personne indépendante, ou montrer que le risque était mal mesuré et faire valider un modèle corrigé. Ces trois mêmes voies sont aussi les voies malhonnêtes, lorsque la
hausse est accordée pour épouser la position, ou que le modèle est modifié pour s'adapter à la limite.

\omcode{code/firm/escalation/firm_escalation.py}{28}{53}{Les règles d'escalade sous forme de données~: gravité selon la taille et la durée, et qui doit être informé
dans combien de jours ouvrés.}\label{lst:fm:the-risk-management-function:rules}

\begin{omfigure}
\begin{tikzpicture}
\begin{axis}[width=10.5cm,height=5.2cm,xlabel={\small mois},ylabel={\small milliards de dollars},xmin=0,xmax=15,ymin=0,ymax=5,
  xtick={0,4,7,10,13,15},xticklabels={nov. 06,mars 07,juin 07,sept. 07,déc. 07,févr. 08},
  tick label style={font=\footnotesize},grid=major,grid style={gray!20},table/col sep=comma,
  legend cell align=left,legend style={font=\scriptsize,at={(0.5,-0.3)},anchor=north,/tikz/every even column/.append style={column sep=8pt}},legend columns=2]
\addplot[omDef,very thick,const plot] table[x=month,y=limit]{figdata/desk/12-the-risk-management-function/lehman.csv};
\addplot[omProp,thick,dashed] coordinates {(13,2.5) (15,2.5)};
\node[font=\scriptsize,align=right,anchor=south east] at (axis cs:12.8,4.05) {approuvée en janv.\ 2008,\\datée rétroactivement au 3 déc.\ 2007};
\node[font=\scriptsize,align=left,anchor=north west] at (axis cs:1.2,3.2) {le directeur des risques a plaidé\\pour 2.6--2.7~milliards de dollars};
\legend{limite d'appétit pour le risque à l'échelle de la société,limite 2008 calculée selon la méthode de 2007}
\end{axis}
\end{tikzpicture}

\omcaption{La limite d'\omterm{def:fm:the-head-of-desks-job:appetite}{appétit pour le risque} à l'échelle de la société dans le rapport de l'examinateur de Lehman~: 2.3~milliards de dollars, portée à 3.3~milliards de dollars pour l'exercice 2007, à
3.5~milliards de dollars le 7 septembre 2007 (la société dépassait la nouvelle limite chaque jour ouvré de septembre sauf un), et à 4.0~milliards de dollars par une approbation du 14
janvier 2008 datée rétroactivement au 3 décembre 2007~; les conseillers de l'examinateur évaluent la limite 2008, calculée selon la méthode de l'année précédente, à environ 2.5~milliards de dollars. Source~:
le registre de référence du chapitre.}\label{fig:fm:the-risk-management-function:lehman}
\end{omfigure}

Les archives publiques (la \cref{fig:fm:the-risk-management-function:lehman}) montrent chacune des sorties malhonnêtes. Une limite relevée alors que la société dépassait
la limite~; une hausse datée rétroactivement pour que le dépassement disparaisse du registre~; une limite recalculée avec de nouvelles hypothèses qui donnait un chiffre plus élevé~;
et, avant celles-ci, une position importante exclue du calcul de l'\omterm{def:fm:limits-and-capital-allocation:breach}{utilisation de la limite} de mai à août 2007, ce qui maintenait cette utilisation sous la limite jusqu'à
son intégration. Le livre~6 raconte un cas postérieur où un modèle de risque a été modifié pendant qu'une limite était dépassée.
Le registre d'escalade doit être construit de sorte que de tels schémas soient visibles~: limites et modèles conservés à leur date, tels qu'ils étaient connus chaque jour, et chaque changement
consigné avec le jour de son approbation ainsi que celui de sa prise d'effet.

\begin{proposition}[Attendre le comité]\label{prop:fm:the-risk-management-function:wait}
Si le \omterm{def:fm:the-risk-management-function:cro}{comité des risques} se réunit tous les $m$-ièmes jours ouvrés et qu'un dépassement survient un jour uniformément réparti sur le cycle, l'attente
jusqu'à la réunion suivante après le jour d'ouverture suit une loi uniforme sur $\{1,\dots,m\}$, d'espérance $(m+1)/2$ jours ouvrés.
\end{proposition}

\begin{proof}
La réunion suivante après le jour $t$ est le jour $t+1+((-(t+1))\bmod m)$~; lorsque $t$ parcourt un cycle, l'attente $1+((-(t+1))\bmod m)$ prend chaque valeur
$1,\dots,m$ une fois.
\end{proof}

Avec des réunions hebdomadaires, l'attente est en moyenne de trois jours ouvrés. Une procédure qui fait attendre le desk avec son dépassement encore ouvert allonge chaque dépassement
qui est porté devant le comité~; une procédure qui oblige le desk à réduire sa position pendant qu'il attend ne l'allonge pas.

\section{Tutoriel~: une année de dépassements}\label{tut:fm:the-risk-management-function:tut}

\begin{tutorial}
\textbf{Objectif.} Faire tourner les règles d'escalade sur une année d'\omterm{def:fm:limits-and-capital-allocation:breach}{utilisation de la limite} pour vingt desks, mesurer comment les dépassements sont résolus et combien de temps ils durent,
et repérer le schéma des archives publiques. \textbf{État final~:} un tableau et deux graphiques sur la façon dont les dépassements se terminent (le \cref{tab:fm:the-risk-management-function:regimes},
la \cref{fig:fm:the-risk-management-function:regimes} et la \cref{fig:fm:the-risk-management-function:window}).
\begin{tutsteps}
\item \textbf{L'année.} \texttt{fm\_risk.simulate(regime)} tire la valeur en risque de chaque desk par rapport à une \omterm{def:fm:limits-and-capital-allocation:breach}{limite ferme} de 100 (une marche aléatoire persistante
autour de 75~\% d'utilisation). Lorsqu'un dépassement de la \omterm{def:fm:limits-and-capital-allocation:breach}{limite ferme} survient, le desk demande un relèvement temporaire (probabilité 0.45), change son modèle de risque (0.08~; le
risque mesuré baisse d'un cinquième) ou réduit sa position.
\item \textbf{Trois régimes.} Le responsable de desk approuve les relèvements le lendemain~; ou le \omterm{def:fm:the-risk-management-function:cro}{comité des risques}, qui se réunit tous les cinq jours, les approuve avec
une probabilité de 0.6 et valide les changements de modèle lors de sa réunion~; ou le même comité, le desk étant tenu de réduire sa position pendant qu'il attend.
\item \textbf{Le registre.} \texttt{firm.escalation.register} repère chaque dépassement, lui attribue un niveau de gravité, classe sa clôture et signale les limites relevées ou
les modèles modifiés pendant un dépassement et les relèvements répétés (le \cref{lst:fm:the-risk-management-function:register}).
\item \textbf{La comparaison.} \texttt{fm\_risk.compare()} regroupe dix années simulées pour chaque régime.
\end{tutsteps}
\end{tutorial}

\omcode{code/firm/escalation/firm_escalation.py}{77}{111}{Le registre des dépassements~: chaque série de jours au-dessus de la limite ferme, notée, classée selon
la manière dont il a été clos et contrôlé pour repérer les signaux d'alerte.}\label{lst:fm:the-risk-management-function:register}

\begin{table}[ht]
\centering\small
\begin{tabular}{>{\raggedright\arraybackslash}p{5.0cm}>{\centering\arraybackslash}p{2.8cm}>{\centering\arraybackslash}p{2.8cm}>{\centering\arraybackslash}p{2.8cm}}
\toprule
 & responsable de desk & \omterm{def:fm:the-risk-management-function:cro}{comité des risques} & comité, \\
 & approuve & approuve & réduction pendant l'attente \\
\midrule
dépassements par an (20 desks) & 27.8 & 26.4 & 27.8 \\
clos par une réduction de position (\%) & 50.0 & 78.4 & 87.4 \\
clos par un relèvement de limite (\%) & 45.3 & 14.4 & 5.4 \\
clos par un changement de modèle (\%) & 3.2 & 6.1 & 6.5 \\
durée moyenne (jours) & 1.15 & 1.64 & 1.25 \\
jours de dépassement par an & 31.8 & 43.5 & 34.6 \\
gravité 2 ou 3 (dix ans) & 32 & 57 & 25 \\
jours-desk au-dessus de la limite initiale (\%) & 2.44 & 1.33 & 0.74 \\
dépassements signalés (\%) & 54.7 & 23.1 & 11.9 \\
\bottomrule
\end{tabular}
\caption{Dix années simulées de dépassements sur vingt desks selon trois régimes d'approbation (0.7--1.4~\% des dépassements restent ouverts à la fin d'une année).
Données~: \texttt{fm\_risk.compare}.}\label{tab:fm:the-risk-management-function:regimes}
\end{table}

La plupart des dépassements sont petits et brefs~: la durée médiane est d'un jour dans chaque régime, car l'utilisation qui dépasse une limite redescend souvent.
Les régimes diffèrent par la façon dont se terminent les autres (le \cref{tab:fm:the-risk-management-function:regimes}). Lorsque le responsable de desk approuve les relèvements, 45~\% des
dépassements sont clos par un relèvement de la limite~; dans ce régime ils sont les plus brefs, et les desks passent 2.44~\% de leurs jours au-dessus de la limite qu'ils avaient en début
d'année. Exiger l'approbation du comité ramène les relèvements à 14~\% et le temps au-dessus de la limite initiale à 1.33~\%, mais les dépassements durent plus longtemps pendant que
les desks attendent une réunion (la \cref{prop:fm:the-risk-management-function:wait})~: 43.5 jours de dépassement par an contre 31.8, et davantage d'entre eux atteignent
un niveau de gravité 2 ou 3. Obliger les desks à réduire leur position pendant l'attente préserve le contrôle du comité et supprime le délai~: 5~\% des dépassements se terminent par un relèvement,
la durée moyenne est de 1.25 jour, et les desks passent 0.74~\% de leurs jours au-dessus de leurs limites initiales.

\begin{omfigure}
\begin{tikzpicture}
\begin{axis}[width=10.5cm,height=4.6cm,xbar stacked,bar width=12pt,xmin=0,xmax=100,xlabel={\small part des dépassements (\%)},
  ytick={0,1,2},yticklabels={responsable de desk approuve,comité des risques approuve,{comité, réduction pendant l'attente}},y dir=reverse,ymin=-0.6,ymax=2.6,
  tick label style={font=\footnotesize},grid=major,grid style={gray!20},table/col sep=comma,
  legend cell align=left,legend style={font=\scriptsize,at={(0.5,-0.35)},anchor=north,/tikz/every even column/.append style={column sep=8pt}},legend columns=4]
\addplot[fill=omDef!55,draw=omDef] table[y=pos,x=cut]{figdata/desk/12-the-risk-management-function/regimes.csv};
\addplot[fill=omProp!55,draw=omProp] table[y=pos,x=increase]{figdata/desk/12-the-risk-management-function/regimes.csv};
\addplot[fill=omThm!50,draw=omThm] table[y=pos,x=model]{figdata/desk/12-the-risk-management-function/regimes.csv};
\addplot[fill=gray!40,draw=gray] table[y=pos,x=open]{figdata/desk/12-the-risk-management-function/regimes.csv};
\legend{réduction de position,relèvement de limite,changement de modèle,encore ouvert}
\end{axis}
\end{tikzpicture}

\omcaption{Comment se terminent les dépassements selon trois régimes d'approbation~: dix années simulées de vingt desks chacun. Retirer au desk l'approbation des relèvements
déplace les dépassements de la limite vers la position. Données~: \texttt{fm\_risk.compare}.}\label{fig:fm:the-risk-management-function:regimes}
\end{omfigure}

\begin{omfigure}
\begin{tikzpicture}
\begin{axis}[width=10.5cm,height=5cm,xlabel={\small jours ouvrés depuis le changement de modèle},ylabel={\small part de la limite initiale},xmin=-15,xmax=24,
  ymin=0.6,ymax=1.25,tick label style={font=\footnotesize},grid=major,grid style={gray!20},table/col sep=comma,
  legend cell align=left,legend style={font=\scriptsize,at={(0.5,-0.3)},anchor=north,/tikz/every even column/.append style={column sep=8pt}},legend columns=2]
\addplot[omDef,thick,mark=*,mark size=1pt] table[x=day,y=utilisation]{figdata/desk/12-the-risk-management-function/window.csv};
\addplot[omProp,thick,const plot] table[x=day,y=limit]{figdata/desk/12-the-risk-management-function/window.csv};
\draw[gray,dotted] (axis cs:0,0.6) -- (axis cs:0,1.25);
\legend{valeur en risque mesurée,limite ferme (avec relèvements temporaires)}
\end{axis}
\end{tikzpicture}

\omcaption{Un desk autour du premier dépassement de l'année simulée, clos par un changement de modèle~: un relèvement temporaire à 110, une interruption d'un jour à son échéance et un nouveau
relèvement, puis un dépassement de 7.7~\% clos deux jours plus tard par un changement de modèle qui réduit la valeur en risque mesurée d'un cinquième, sans modification de la position. Le
registre signale à la fois les relèvements répétés et le changement de modèle. Données~: \texttt{fm\_risk.model\_change\_window}.}\label{fig:fm:the-risk-management-function:window}
\end{omfigure}

\textbf{Pour aller plus loin.} Notez chaque relèvement avec son jour d'approbation et vérifiez-le avec \texttt{firm.escalation.backdated}~; donnez aux desks une
position qui est exclue de l'utilisation pendant trois mois et voyez quelle métrique du registre la signale en premier.

\section{Ce que recherche le responsable des risques}

Un registre qui se contente de compter les dépassements manque les défaillances du dossier public, qui relevaient de schémas plutôt que d'événements. Le responsable des risques recherche~:
\begin{itemize}
\item des limites relevées, ou des modèles modifiés, alors qu'un dépassement est ouvert, et des relèvements temporaires que l'on renouvelle au lieu de les laisser expirer~;
\item des changements qui prennent effet avant d'être approuvés~;
\item des positions hors de la mesure~: expositions exclues du calcul de l'utilisation ou des tests de résistance parce qu'elles sont nouvelles, peu liquides ou
difficiles à modéliser (l'examinateur a constaté que les tests de résistance de l'établissement excluaient ses investissements immobiliers commerciaux et de capital-investissement)~;
\item des limites recalculées avec de nouvelles hypothèses qui donnent un chiffre plus élevé au moment où un chiffre plus élevé est nécessaire~;
\item des \omterm{def:fm:the-risk-management-function:escalation}{dérogations de risque} sans responsable ni échéance, et un comité qui approuve tout ce qu'on lui soumet.
\end{itemize}

\begin{definition}[Culture du risque]\label{def:fm:the-risk-management-function:culture}
La \emph{culture du risque}\index{culture du risque} est l'ensemble des normes, des attitudes et des comportements d'un établissement qui déterminent la manière dont le risque est pris et contrôlé en
pratique~: si les limites sont traitées comme contraignantes, si les mauvaises nouvelles remontent vite, et si ceux qui s'opposent sont entendus.
\end{definition}

La culture ne se lit pas directement dans un registre, mais ses traces, elles, s'y lisent~: la part des dépassements clos par des relèvements, la part des dépassements signalés, le temps
écoulé entre un dépassement et le moment où le comité en est informé, le nombre de dérogations dépassant leur échéance. L'examinateur a constaté que la direction avait traité la
limite globale comme une simple ligne directrice «~soft~», tout en la présentant à son régulateur et à son conseil comme une contrainte réelle. Sur la question
que devait trancher l'examinateur, le rapport n'a pas trouvé d'éléments suffisants pour établir que la conduite des hauts dirigeants dans la gestion du risque sortait de la règle dite business judgement rule ou
était téméraire ou irrationnelle~: les décisions qui se sont révélées malavisées étaient, juridiquement, des décisions que la direction avait le droit de prendre. La leçon pour une fonction de gestion des risques
est que son autorité doit être inscrite dans la procédure, car elle ne peut pas compter sur le droit pour la lui donner après coup.

\begin{remark}[Les limites du chapitre~7 et la procédure du présent chapitre]\label{rem:fm:the-risk-management-function:ch7}
Le chapitre~7 dimensionnait les limites et tarifait le capital comme si chaque limite tenait. Le registre de ce chapitre permet à un établissement de savoir si elles tiennent~: une limite qui est
relevée chaque fois qu'elle devient contraignante n'est pas une limite, et le capital qu'elle devait protéger ne l'est pas.
\end{remark}

\section{Construction~: le registre d'escalade}\label{bld:fm:the-risk-management-function:escalation}

\begin{build}
\bfield{Rôle} Le registre des dépassements et le flux d'escalade sous forme de données~: classement, notification, suivi de la résolution, métriques d'audit et signaux d'alerte.
\bfield{Interface} \texttt{firm.escalation}~: \texttt{Rules}, \texttt{severity}, \texttt{notify}, \texttt{hard\_path} (à partir d'un
nœud \texttt{firm.limitalloc} et de relèvements datés), \texttt{register}, \texttt{stats}, \texttt{flagged}, \texttt{backdated}.
\bfield{Règles} Un dépassement est une suite de jours au-dessus de la \omterm{def:fm:limits-and-capital-allocation:breach}{limite ferme} effective~; il se clôt par un relèvement de limite, un changement de modèle ou une réduction de position, dans cet
ordre de priorité~; les signalements sont enregistrés, jamais effacés~; les limites et les modèles sont pris à la date considérée.
\bfield{Tests d'acceptation} \texttt{code/firm/escalation/tests/}~: la grille de gravité~; un dépassement de chaque type de résolution sur des données construites à la main~; des relèvements
répétés~; un changement antidaté.
\bfield{Pour aller plus loin} Délais de notification vérifiés par un journal de qui a été informé et quand~; dérogations assorties de responsables et de dates d'échéance~; un rapport quotidien pour le
\omterm{def:fm:the-risk-management-function:cro}{comité des risques}.
\end{build}

\begin{omsources}
\item Report of Anton R. Valukas, Examiner, In re Lehman Brothers Holdings Inc., No.~08-13555 (Bankr. S.D.N.Y.), 11 mars 2010, volume~1.
\item 12 CFR 252.31 and 252.33 (Regulation YY).
\item Senior Supervisors Group, \emph{Observations on Risk Management Practices during the Recent Market Turbulence}, 6 mars 2008.
\end{omsources}

\section{Exercices}

\begin{exercise}[$\star$]\label{exo:fm:the-risk-management-function:1}
Classer selon les règles du chapitre un dépassement de 3~\% durant un jour, un de 10~\% durant un jour et un de 3~\% durant six jours, et dire qui est informé et
dans quels délais.
\end{exercise}

\begin{exercise}[$\star$]\label{exo:fm:the-risk-management-function:2}
De quel pourcentage la limite globale de la \cref{fig:fm:the-risk-management-function:lehman} a-t-elle augmenté de novembre 2006 à décembre 2007~? De combien
la limite de 2008 a-t-elle dépassé celle que donnait la méthode de l'année précédente~?
\end{exercise}

\begin{exercise}[$\star$]\label{exo:fm:the-risk-management-function:3}
Un \omterm{def:fm:the-risk-management-function:cro}{comité des risques} se réunit tous les cinquièmes jours ouvrés. Combien de temps, en moyenne, un desk attend-il ce comité après l'ouverture d'un dépassement~? Et avec des réunions tous les
dix jours~?
\end{exercise}

\begin{exercise}[$\star\star$]\label{exo:fm:the-risk-management-function:4}
Pourquoi un relèvement de limite antidaté est-il un signal d'alerte même lorsque le relèvement lui-même est justifié~?
\end{exercise}

\begin{exercise}[$\star\star$]\label{exo:fm:the-risk-management-function:5}
Dans le \cref{tab:fm:the-risk-management-function:regimes}, pourquoi les dépassements durent-ils plus longtemps lorsque le comité approuve les relèvements, et pourquoi exiger
que le desk réduise sa position pendant qu'il attend fait-il redescendre la durée~?
\end{exercise}

\begin{exercise}[$\star\star$]\label{exo:fm:the-risk-management-function:6}
Expliquer pourquoi un responsable de desk qui approuve les relèvements pour son propre desk respecte le \omterm{def:fm:the-risk-management-function:foureyes}{principe des quatre yeux} sur le papier mais pas dans les faits.
\end{exercise}

\begin{exercise}[$\star\star\star$]\label{exo:fm:the-risk-management-function:7}
\emph{Programmation.} Lancer \texttt{fm\_risk.run} pour le régime des desks et lister les dépassements du desk 19 avant son changement de modèle. Lesquels sont signalés, et pourquoi~?
\end{exercise}

\begin{exercise}[$\star\star\star$]\label{exo:fm:the-risk-management-function:8}
\emph{Trouver la faille.} «~Notre fonction de gestion des risques est indépendante~: le \omterm{def:fm:the-risk-management-function:cro}{directeur des risques} (CRO) rend compte au responsable du trading, qui comprend les risques mieux
que quiconque.~»
\end{exercise}

\section{Problème~: Relever pour tenir}

\begin{problem}[{Problème du week-end --- relever pour tenir}]\label{pb:fm:the-risk-management-function:1}
Un nouveau \omterm{def:fm:the-risk-management-function:cro}{directeur des risques} (CRO) hérite d'un établissement où les responsables de desk approuvent leurs propres relèvements temporaires, et doit proposer une
nouvelle \omterm{def:fm:the-risk-management-function:escalation}{procédure d'escalade} au \omterm{def:fm:the-risk-management-function:cro}{comité des risques}.

\medskip\noindent\textbf{Partie I --- La fonction.}
\begin{enumerate}
\item Définir le \omterm{def:fm:the-risk-management-function:cro}{directeur des risques} (CRO) et le \omterm{def:fm:the-risk-management-function:cro}{comité des risques}.
\item Quelles sont les trois conditions qui rendent le \omterm{def:fm:the-risk-management-function:cro}{directeur des risques} indépendant dans les faits~?
\item Qu'est-ce que le Senior Supervisors Group a identifié comme trait distinctif des établissements qui ont évité des pertes importantes en 2007~?
\item Définir le \omterm{def:fm:the-risk-management-function:foureyes}{principe des quatre yeux} et donner deux décisions qu'il doit couvrir.
\end{enumerate}

\medskip\noindent\textbf{Partie II --- Le dossier public.}
\begin{enumerate}[resume]
\item Donner l'évolution de la limite globale dans le rapport de l'examinateur, avec les dates.
\item Quel effet l'antidatage de janvier 2008 a-t-il eu sur le registre des dépassements~?
\item Qu'ont laissé de côté les tests de résistance, et quelle position a été exclue du calcul de l'utilisation~?
\item Que l'examinateur a-t-il conclu au sujet de la business judgement rule, et que cela signifie-t-il pour une fonction de gestion des risques~?
\end{enumerate}

\medskip\noindent\textbf{Partie III --- Le registre.}
\begin{enumerate}[resume]
\item Définir une \omterm{def:fm:the-risk-management-function:escalation}{procédure d'escalade} et une \omterm{def:fm:the-risk-management-function:escalation}{dérogation de risque}.
\item Énoncer les règles de gravité du chapitre et classer trois dépassements de votre choix.
\item Énoncer et démontrer la \cref{prop:fm:the-risk-management-function:wait}.
\item Donner la part des dépassements clos par des hausses dans les trois régimes.
\item Donner la durée moyenne, le nombre de jours de dépassement par an et le temps passé au-dessus de la limite initiale dans chacun.
\item Que signale le registre, et quelle part des dépassements est signalée dans chaque régime~?
\end{enumerate}

\medskip\noindent\textbf{Partie IV --- La proposition.}
\begin{enumerate}[resume]
\item Pourquoi la durée médiane est-elle la même dans chaque régime, et quelle statistique montre la différence~?
\item Décrire l'épisode de changement de modèle de la \cref{fig:fm:the-risk-management-function:window}.
\item Quels relevés ponctuels le registre doit-il conserver~?
\item Définir la \omterm{def:fm:the-risk-management-function:culture}{culture du risque} et nommer trois métriques du registre qui en témoignent.
\item Énoncer le \emph{résultat nommé}~: la part des dépassements résolus par la hausse de la limite, le temps médian de résolution, et la façon dont les deux évoluent lorsque
les hausses requièrent l'approbation du \omterm{def:fm:the-risk-management-function:cro}{comité des risques}.
\item En deux phrases, rédiger la nouvelle \omterm{def:fm:the-risk-management-function:escalation}{procédure d'escalade}.
\end{enumerate}
\end{problem}

\section{Questions d'entretien}

\begin{interviewq}[$\star$ \iqroles{risque}]\label{iq:fm:the-risk-management-function:1}
À qui un \omterm{def:fm:the-risk-management-function:cro}{directeur des risques} devrait-il rendre compte, et pourquoi~?
\end{interviewq}

\begin{interviewq}[$\star$ \iqroles{trader, risque}]\label{iq:fm:the-risk-management-function:2}
Votre desk dépasse sa limite de VaR de 3~\% à la clôture. Que se passe-t-il ensuite~?
\end{interviewq}

\begin{interviewq}[$\star\star$ \iqroles{risque}]\label{iq:fm:the-risk-management-function:3}
Un desk demande une troisième hausse temporaire en deux mois. Que faites-vous~?
\end{interviewq}

\begin{interviewq}[$\star\star$ \iqroles{risque, chercheur}]\label{iq:fm:the-risk-management-function:4}
Un desk propose un nouveau modèle de VaR qui abaisse son risque mesuré de 20~\% alors qu'il est en dépassement. Comment traitez-vous cela~?
\end{interviewq}

\begin{interviewq}[$\star\star$ \iqroles{risque}]\label{iq:fm:the-risk-management-function:5}
Que chercheriez-vous dans le registre des dépassements d'une société pour juger de sa \omterm{def:fm:the-risk-management-function:culture}{culture du risque}~?
\end{interviewq}

\begin{interviewq}[$\star\star\star$ \iqroles{risque, développeur}]\label{iq:fm:the-risk-management-function:6}
Concevez le modèle de données d'un registre des dépassements qui ne puisse pas être réécrit après coup.
\end{interviewq}
